\section{Resultados}
\label{sec:resultados}

\subsection{Presupuesto de tiempo real: latencias medidas}
\label{sec:res_latencias}

La Tabla~\ref{tab:latencias} resume las latencias medidas.
El presupuesto de diseño (pausa$\to$enrutamiento $< 1$ s; cuadro$\to$etiqueta
$< 250$ ms) se cumple con margen.

\begin{table}[H]
\centering
\caption{Latencias medidas en el equipo de desarrollo (RTX 2050; ResNet18 en CPU).}
\label{tab:latencias}
\begin{tabular}{llr}
\toprule
Canal & Operación & Latencia \\
\midrule
Voz    & Transcripción Whisper \emph{small} (GPU, frase corta) & 336 ms \\
Voz    & Enrutamiento de texto (8 agentes, broadcast gateado)  & 302 ms \\
Voz    & \textbf{Extremo a extremo} (fin de habla $\to$ decisión) & $\approx$ 0.6 s \\
\midrule
Visión & Análisis de cuadro (encode $+$ directorio visual)     & 54 ms \\
Visión & Interfaz con cámara activa                            & 26 fps \\
\midrule
Evocación & Recall SS $+$ prototipo $+$ distancias $+$ decodificación
          & $\approx$ 2 s (asíncrona) \\
\bottomrule
\end{tabular}
\end{table}

La estructura de estas cifras confirma la hipótesis de viabilidad: el enrutamiento es
proyección matricial y cuesta milisegundos por agente; los cientos de milisegundos del
canal de voz pertenecen al traductor de frontera (Whisper), no a las memorias.
La única operación intrínsecamente lenta ---la recuperación estocástica--- no
participa del lazo de tiempo real: corre asíncrona y su resultado se incorpora al
panel al estar listo, sin interrumpir video ni escucha.

\subsection{Cadena de voz extremo a extremo}

La prueba inyectada reproduce el caso de uso completo sin micrófono: un WAV con la
frase ``\emph{Farm animal that gives milk}'' atraviesa VAD, transcripción,
enrutamiento y evocación.
El resultado: transcripción correcta, asignación a \textit{cow} en modo temprano
(con el registro correspondiente en el directorio de sesión), y en modo maduro la
redirección punto a punto \textit{apple} $\to$ \textit{cow}
(\texttt{routed=True}, 56 ms) con entrada \textit{apple}, y la permanencia
\textit{cow} $\to$ \textit{cow} (``se queda: es el experto'') con entrada
\textit{cow}.
La evocación asíncrona produjo la imagen del recuerdo $\sim$2 s después del
enrutamiento (token \texttt{farm}, peso 9142) sin congelar la interfaz
(Figura~\ref{fig:wavtest}).

\begin{figure}[H]
\centering
\includegraphics[width=0.35\textwidth]{wavtest_evocacion}
\caption{Imagen evocada por la cadena de voz inyectada (``Farm animal that gives
milk'' $\to$ \textit{cow}): reconstrucción del decoder a partir del patrón recuperado
por la memoria hetero-asociativa, no una fotografía del conjunto de datos.}
\label{fig:wavtest}
\end{figure}

En el micrófono real, la sonda de arranque detectó el caso límite relevante del
sistema operativo: con el permiso de privacidad de Windows desactivado la señal llega
en silencio absoluto (RMS $= 0$) y el sistema lo reporta con instrucciones; con el
permiso activo el arreglo de micrófonos entrega señal normal (RMS $\approx 0.0013$ en
reposo).

\subsection{Canal de visión: sintético y físico}

Con cuadros sintéticos (imagen ETH-80 llenando el recuadro al 55\%), el canal enruta
correctamente en $\sim$98 ms y rechaza el cuadro vacío mostrando el candidato más
cercano.
Con cuadros ETH-80 reales presentados sin degradación, \textit{car} enruta a
\textit{car} y las clases con menor margen caen en rechazo ---consistente con el
comportamiento del hemisferio visual del experimento base (75.0\% de enrutamiento con
cero falsos aceptados y un 25\% de rechazo residual): el laboratorio en vivo hereda
los límites del experimento, no agrega los suyos.

La prueba física arrojó el diagnóstico central de esta sección: la \textbf{pantalla de
teléfono no es un estímulo viable} (brillo, moiré y encuadre colapsan el enrutamiento,
tal como lo anticipó la simulación), mientras que la \textbf{hoja impresa en papel
mate} sí lo es, siempre que el recuadro de análisis ajustable permita llenarla a la
distancia de enfoque de la cámara.
La primera sesión con impresiones dejó solo a \textit{cow} como clase fiable; el
análisis de los cuadros capturados mostró el corrimiento de color de la cadena
impresora$+$cámara como causa, lo que motivó la política de la
Sección~\ref{sec:arquitectura} y el banco siguiente.

\subsubsection*{Política de color: validación}

\begin{table}[H]
\centering
\caption{Enrutamiento correcto sobre el banco de 24 imágenes (3 por clase) bajo
condiciones simuladas de cámara. La política aplica el rescate
\emph{solo} tras un rechazo del intento crudo.}
\label{tab:colorfix}
\begin{tabular}{lccc}
\toprule
Condición & Sin rescate & Con política completa & El rescate actuó en \\
\midrule
Limpia             & 22/24 & 22/24 & 1/24 \\
Degradación media  & ---   & 18/24 & --- \\
Degradación fuerte & $\sim$0/24 & 15/24 & mayoría de los aciertos \\
\bottomrule
\end{tabular}
\end{table}

Dos lecturas de la Tabla~\ref{tab:colorfix}.
Primero, el rescate condicionado \emph{no daña} las entradas limpias (22/24 en ambos
casos): al correr solo tras un rechazo, una corrección innecesaria nunca toca una
entrada ya aceptada.
Segundo, bajo degradación fuerte el rescate recupera el canal visual desde el colapso
($\sim$0) hasta 15/24.
El hallazgo negativo también se reporta: un detector previo por distancia de momentos
es inviable porque las distribuciones limpia/degradada se traslapan (medianas 24.2
vs.\ 27.1) --- de ahí que la política sea ``reintentar tras rechazo'' y no
``detectar y corregir''.

\subsection{Pista faltante: prototipo, RS, ST y SS}
\label{sec:res_pista}

La Tabla~\ref{tab:pista} reporta la caracterización distribucional de los cuatro
métodos de lectura del plano proyectado con la distancia de
retro-proyección~\eqref{eq:test} como métrica común: 30 corridas con semillas fijas
por método y pista (el prototipo es determinista).

\begin{table}[H]
\centering
\caption{Distancia de retro-proyección (menor = mejor) de las cuatro lecturas del
plano proyectado, sobre la misma proyección por pista: media $\pm$ desviación
estándar sobre 30 corridas con semillas fijas, y fracción de corridas que alcanzan
distancia exactamente cero. El prototipo es determinista (un valor por pista).}
\label{tab:pista}
\begin{tabular}{llrrrrrr}
\toprule
 & & & \multicolumn{3}{c}{media $\pm$ $\sigma$ (30 corridas)}
 & \multicolumn{2}{c}{corridas con $d = 0$} \\
\cmidrule(lr){4-6}\cmidrule(lr){7-8}
Pista & Agente & Prototipo & RS & ST & SS & ST & SS \\
\midrule
\texttt{pear}   & \textit{pear}  & 9.02 & $7.82 \pm 1.67$ & $0.36 \pm 1.12$ & $0.06 \pm 0.32$ & 27/30 & 28/30 \\
\texttt{car}    & \textit{car}   & 7.97 & $4.37 \pm 1.45$ & $0.03 \pm 0.16$ & $0.01 \pm 0.03$ & 29/30 & 29/30 \\
\texttt{green}  & \textit{apple} & 7.00 & $4.00 \pm 1.25$ & $0.00 \pm 0.00$ & $0.00 \pm 0.00$ & 30/30 & 30/30 \\
\texttt{animal} & \textit{horse} & 7.01 & $5.02 \pm 1.02$ & $1.17 \pm 1.08$ & $0.19 \pm 0.23$ & 13/30 & 13/30 \\
\bottomrule
\end{tabular}
\end{table}

Los costos medios por corrida: RS $\approx$ 0.02 s (una muestra), ST 0.9--1.4 s
(127 muestras con prueba), SS 1.9--5.8 s según la pista (ST más descenso local; la
búsqueda se alarga justo en las pistas difíciles).
La jerarquía del artículo \cite{morales2025missing} se reproduce en este dominio:
$\mathrm{SS} \le \mathrm{ST} < \mathrm{RS} < \text{prototipo}$ en las cuatro pistas,
con desigualdad estricta $\mathrm{SS} < \mathrm{ST}$ en las pistas difíciles
(\texttt{pear}, \texttt{car}, \texttt{animal}) y empate en el techo ($0.00$) en
\texttt{green}.
Tres observaciones:

\begin{enumerate}
    \item \textbf{El descenso local de SS aporta valor exactamente donde la pista es
    difícil.} En las pistas fáciles ST y SS son casi indistinguibles
    (\texttt{green}: ambos $0.00$; \texttt{car}: $0.03$ vs.\ $0.01$), pero en las
    difíciles SS reduce la media de ST a una fracción: en \texttt{animal}
    ---una pista semánticamente compartida entre las clases de animales---
    $1.17 \pm 1.08 \to 0.19 \pm 0.23$, y en \texttt{pear}
    $0.36 \pm 1.12 \to 0.06 \pm 0.32$.
    Cuando el muestreo no alcanza el óptimo (en \texttt{animal} solo 13 de 30
    corridas llegan a cero), la búsqueda acerca las corridas restantes.
    \item \textbf{RS es errático y rara vez acierta}: medias de 4.0 a 7.8 con
    desviaciones de $\pm 1.0$ a $\pm 1.7$, y distancia cero en apenas 2 de 120
    corridas totales --- el comportamiento esperado de la línea base sin prueba.
    \item \textbf{El prototipo es consistentemente el peor} (7.0--9.0): queda por
    encima incluso de la \emph{media} de RS en las cuatro pistas, pese a ser la
    lectura determinista ``más probable'' columna a columna.
    Elegir el máximo de cada característica por separado destruye las correlaciones
    conjuntas que el registro sí almacena: el resultado es la ``plastilina''
    ---una superposición estadística--- y no un objeto que la memoria haya contenido.
\end{enumerate}

\begin{figure}[H]
\centering
\includegraphics[width=0.9\textwidth]{panel_prototipo_definido}
\caption{Panel de la ventana en tiempo real tras la frase ``animal with a mane''
$\to$ \textit{horse}: arriba el prototipo emergente del plano proyectado
(\emph{prototype, raw plane}, $d = 7.0$), visiblemente difuso; abajo el objeto
definido del recall oficial SS (\emph{memory of horse, defined}, $d = 0.0$).
A la izquierda, las barras del enrutamiento por agente y el directorio de sesión.}
\label{fig:panel}
\end{figure}

El contraste visual de la Figura~\ref{fig:panel} es el resultado cualitativo central:
el prototipo decodifica a una mancha reconocible pero desdibujada (la superposición),
mientras que el objeto definido decodifica a una figura nítida.
La misma comparación, extendida a los cuatro métodos y al mapa de calor del plano
proyectado, está disponible en la etapa de recall de la aplicación de trazado del
experimento.
